\section{System Overview}
\label{sec:system_overview}

\begin{figure*}[t]
    \centering
    \includegraphics[width=\textwidth]{figures/1_system_overview_v7_compressed.pdf}
    \vspace{-6mm}
    \caption{
    Our system takes the scene description and human-level instruction as input and uses a high-level planner to obtain the scene map and a detailed execution plan. The low-level motion generator then generates synchronized object motion, full-body human motion, and finger motion. Finally, the physics tracker uses RL to track the generated motion, producing physically plausible motion.}
    % \Description{}
    \vspace{-4mm}
    \label{fig:system_overview}
\end{figure*}
Given a human-level instruction and a scene description, we develop a system that generates synchronized human motion and object motion to accomplish the task. 
%Our system differs from prior work in human-scene and human-object interaction by not only modeling the full-body dynamics of human motion but also focusing on the intricate details of finger movements. This enhancement contributes significantly to the realism of the motions generated.
An overview of our system is depicted in \cref{fig:system_overview}. The system consists of a high-level planner, a low-level motion generator, and a physics tracker. The high-level planner leverages LLMs to reason about instructions and the scene description, generating a \emph{scene map} along with a detailed \emph{execution plan}. The scene map specifies the desired positions and orientations of the objects, while the execution plan outlines the steps to interact with these objects. The low-level motion generator unites conditional diffusion models trained on separate datasets for full body motion and dexterous hand motion, producing synchronous object and human movements.
%It includes an interaction module and a navigation module. The interaction module, introduced in \cref{sec:low_level_interaction}, employs a multi-stage framework to produce human-object interaction motions. The navigation module, described in \cref{sec:low_level_navigation}, generates walking motions between interactions.
The physics tracker then uses RL to track the generated motion in a physics simulator, ensuring physical plausibility.
% Together, these modules operate sequentially to ensure smooth and realistic human-object interactions.






